To estimate structural $\gamma$, I use 3 linear models, increasing in complexity,
each resolving some of the endogeneity and identification issues of the last.
Though the final model I present has the most accurate interpretation of parameters,
earlier models are useful for iteratively understanding the data
and making choices about parameter selection along the way
since they are computationally easier to calculate and
assuming that relationships between these parameters are similar.
Choosing all these parameters once, in a final model, would have required
a large amount of computation via cross-validation,
which this investigations limits on time and computation did not permit.

\begin{table}[htbp]
    \begin{center}
        \begin{threeparttable}
            \caption{\todo{Summary of econometric models}}
            \label{table:model-summary}
            \begin{tabular}{lcllcl}
                \toprule\toprule
                Model
                & Stages                             
                & Identification
                & Endog.
                & Consistent?
                & Decisions
                \\
                \hline
                1. \nameref{ss:model-lmm}
                & 2                                     \\
                2. \nameref{ss:model-dd}
                & 2                                     \\
                3. \nameref{s:model-dynamic-panel}
                & 1                                     \\
                \bottomrule
            \end{tabular}
            \begin{tablenotes}
                \item 'Stages' indicates whether the derived 2-factor TFP value $z_{it}$
                was used as the dependent variable, or the model used $y_{it}$ from the \itx{gva1} variable directly.
            \end{tablenotes}
        \end{threeparttable}
    \end{center}
\end{table}

A naïve OLS would fail strict exogeneity in this setting, as error exogeneity is more limited.
\begin{align}
    \mathbb E[\varepsilon_{it}|x_{i\tau}] &=0\ \forall t\geq \tau
    \\
    \notag
    \delta y_{i,\tau-k}, \delta k_{i,\tau} &\in x_{i,t},\ k\leq p
\end{align}

% Section 4: static %


It is crucial that I have multiple, non-overlapping groups in a single model,
where my target firm $z_{it}$ is a member of all groups simultaneously,
but peers in one group are excluded from all others.
\par
The non-overlapping property prevents multicollinearity.
If I were to use nested, overlapping geographies,
namely keeping all firms within the \itx{ttwa} area in that group average,
including the ones I have just counted in the \itx{pc4} and itx{pc8} groups,
then my group averages mechanically align:
$z_{jt,pc8}$ is embedded within $z_{jt,pc4}$, and both are within $z_{jt,ttwa}$,
making $\beta_g$ parameter estimates unstable.
\par
The single model approach prevents omitted variable bias (OVB).
$\mathbb E[z_{jt,ttwa|g}] \propto  \mathbb E[z_{jt,pc8|g}]$,
so running three models where the former is just my single regressor means
the latent $\beta_pc8 \in u_{it}$ is correlated with both $z_{itg}$ and $z_{jtg}$.
We could reasonably argue that the inner rings are causal,
so this causal factor excluded from the regression influences both
and creates heavily upward bias on my outer-ring estimations.
Indeed, this critique extends within further inner rings.
If I think the causal unit is not peer firms within the \itx{pc8} group at all,
but a further hyper-local effect - perhaps those on the same street or same building,
then this latent factor will correlate with both my outcome and regressor and drive the results.
However, given the small number of peers contained within an \itx{pc8} unit and
the limits of data availability to identify firms at more local granularities,
I proceed assuming the effect from \itx{pc8} peers is homogenous.
\par
Finally, the approach using multiple groups as regressors at once gives identification
by breaking transitivity.
If I were to run the same regression with merely one group, say \itx{pc8},
then my regression would suffer the reflection  problem outlined in \citep{Manski1993}.
As I am estimating the expected value of $z_{it}$ and $z_{jt}$ from the same sample,
they cannot be uniquely identified in reduced form.

\noindent
\textit{Reflection problem with a single regressor, $\boldsymbol\iota$ denotes a column vector of $1$s}

\begin{align*}
    \mathbb E[z_{itg}|g] &= \alpha_i + \gamma_t + \beta_0 + \beta_1 \mathbb{E}[z_{jtg}|g]
    \\
    &= (\alpha_i + \gamma_t + \beta_0)\boldsymbol\iota + \beta_1 \mathbf\Omega_t \mathbb E[\mathbf{z}_{tg}|g]
    \\
    &= (\mathbb I - \beta_1 \mathbf\Omega_t)^{-1} (\alpha_i + \gamma_t + \beta_0)\boldsymbol\iota
\end{align*}
%
This problem is mitigated with having 3 distinct groups to which $z_{it}$ belongs.
The results of the LMM estimation implementing equation \ref{eq:lmm-structural1} are below:

% Section 5 %

\begin{align}
    \label{eq:dyn-panel_structural}
    y_{i,t} &= \alpha_i + \gamma_t + \beta_0
                + \beta_1 k_{i,t}
                + \beta_2 l_{i,t}
                + \delta_0 \sum_{j\neq i} w_{ij,t}z_{j,t}
                + \varepsilon_{i,t}
    \\
    &= \alpha_i + \gamma_t + \beta_0
                + \beta_1 k_{i,t}
                + \beta_2 l_{i,t}
                + \delta_0 \sum_{j\neq i} w_{ij,t} \cdot
                    (\alpha_j + \gamma_t + \beta_0 + \varepsilon_{j,t})
                + \varepsilon_{i,t}
\end{align}

To support this argument, 
such that lagged realisations of the dependent variable are used as additional regressors.
I can implement these lags exploiting the rich time series observations available in Fame.
I employ a dynamic specification to improve the plausibility of $\sum_{j\neq i}w_{ij}z_{j,t}$
as an exogenous regressor so that $\delta$ is consistent of the true $\delta_0$.
to help eliminate 
\par
Autoregressive lags help control for reverse causality: if $y_{i,t-k}$ is causal for $z_{jt}$,
or feasibly $z_{j,t-1}$ and $\mathbb E[z_{j,t-1}, z_{j,t}]\neq0$, then the presence of
$y_{i,t-k}$ as a control helps eliminate that bias.
Similarly, theory offers us no clear view on where there could be simultaneity,
where variation in $\Delta y_{i,t}$ is driven causally by $\Delta y_{i,t-1}$
and that change in $\Delta y_{i,t}$ in turn drives $\Delta z_{j,t}$.
The autoregressive regressor helps control for that.
Additionally, the results below suggest the autoregressive parameter coefficient is low,
with a log elasticity of 0.005 (significant at the 1\% level) in all specifications,
suggesting any concerns about reverse causality or simultaneity driven by an AR process
are likely minimal.
The use of this dynamic regressor improves the plausible exogeneity of $z_{i,j,t}$.
\par
The autoregressive specification requires a return to the one-stage production function
with an outcome of GVA, $y_{it}$. I continue to use the scalar weighted distance TFP metric
as the primary regressor of interest $z_{jt}$. While introducing dynamic lags of 

\begin{align}
    y_{i,t} &= \beta_0 + \alpha_i + \gamma_t
                + \sum_{k=1}^p \phi_k y_{i,t-k}
                + \beta_1 k_{i,t}
                + \beta_2 l_{i,t}
                + \sum_{k=1}^q \delta_q \sum_{j\neq i} w_{i,j,t}z_{j,t}
                + \varepsilon_{i,t}
\end{align}
%
I list an AR(p,q) specification above,
however, my results show that AR(1), or $p=1$, is the true structure.
Appendix I, which increases the order using a nested model approach,
shows that for any specification order above an AR(1)
coefficients become insignificant suggesting multicollinearity \todo{F-test on AR(2)}.
I also confirm this via information criteria tests, \todo{which show}
that lags beyond $t-2$ provide no new information for the variation on $ln(y_{i,t}$).
Noting that the regression coefficient on the $ln(gva_1)_{t-1}$ is significant at the 1\% level
in my results, I can strongly infer that an AR(1) is the true structure.
\par
\todo{(IC test results)}
